\section{Applications and Connections}
\label{sec:applications}

This work does not propose new dynamics, a new interpretation, or a replacement formalism for
quantum theory. It provides a structural diagnostic: it identifies when an operational reduction
that discards causally relevant information cannot support counterfactual prediction under a
given intervention class, and it characterizes the minimal \emph{representational structure}
required to restore such predictive capacity. We now connect the $\Sigma_1/\Sigma_2$
classification to several concrete and active problem areas.

\subsection{Non-Markovian open quantum systems}
\label{sec:oqs}

Open quantum systems (OQS) theory typically seeks an effective evolution equation for a reduced
state $\rho_S(t)$, e.g.\ Lindblad or memory-kernel master equations, under assumptions such as
initial factorization or weak coupling. In our classification, these are $\Sigma_1$ descriptions:
prediction proceeds by applying an algebraically closed update rule on a reduced object.

The structural result of Section~3 identifies a standard mechanism by which $\Sigma_1$
descriptions can fail: whenever the reduction $\Pi$ identifies distinct correlated histories
as the same reduced state and the intervention class includes free evolution (or more generally,
a sufficiently rich set of admissible post-reduction operations), there need not exist a
single-valued rule on $\rho_S$ that supports counterfactual prediction across all admissible
histories.

This reframes persistent difficulties in strongly coupled or initially correlated regimes as a
\emph{structural} limitation of the reduction: the many-to-one nature of $\Pi$ induces a
degeneracy in continuation structure that cannot be repaired by improving the functional form
of a master equation on $\mathcal{R}$ alone. Classic discussions of how initial correlations
obstruct completely positive reduced dynamics can be viewed as early instances of this
unfaithfulness diagnosis \cite{Pechukas1994,Alicki1995}.

\subsection{Quantum control under memory and uncertain preparation}
\label{sec:control}

Quantum control practice often assumes that a reduced state (or an estimated effective process
model) suffices to predict responses to subsequent control pulses. In correlated or drifting
environments, however, identical $\Sigma_1$ reduced descriptions can respond differently to the
same control intervention. In our language, preparation ambiguity corresponds to fiber
degeneracy $\mathcal{F}(r)=\Pi^{-1}(r)$, and control success depends on which fiber element is
realized.

From the $\Sigma_2$ perspective, this suggests two complementary strategies:
\begin{enumerate}
  \item \textbf{Fiber localization:} augment experiments to infer which subset of $\mathcal{F}(r)$
        is realized (e.g.\ by probing correlation-sensitive observables).
  \item \textbf{Fiber-robust control:} design interventions $I$ such that the induced continuation
        set is effectively insensitive across the relevant equivalence classes of $\mathcal{F}(r)$,
        i.e.\ $\mathcal{A}_I(r)$ is approximately constant over the uncertainty region.
\end{enumerate}
This reframes the control design problem from \emph{reach a target state from $r$} to
\emph{achieve the desired outcome uniformly across a fiber or coarse-grained fiber class}.

\subsection{Preparation contextuality as an unfaithful cut}
\label{sec:contextuality}

Preparation contextuality formalizes the phenomenon that distinct preparation procedures can be
operationally equivalent at the level of a state description, yet yield different statistics
under subsequent measurement choices \cite{Spekkens2005}. Let $\Pi$ map preparation histories
to density matrices. Then preparation contextuality asserts the existence of $h_1,h_2$ with
$\Pi(h_1)=\Pi(h_2)=\rho$ but
$\mathrm{Cont}(\Pi(h_1);\mathcal{I})\neq\mathrm{Cont}(\Pi(h_2);\mathcal{I})$
for an intervention class $\mathcal{I}$ that includes measurement operations. This is precisely
the definition of an unfaithful cut.

From the $\Sigma_2$ perspective, the appropriate descriptive object for preparations is not
$\rho$ alone, but the lifted structure
$\Sigma_2(\rho)=(\rho,\mathcal{F}(\rho),\{\mathcal{A}_I(\rho)\}_{I\in\mathcal{I}})$,
which retains the intervention-indexed continuation geometry required for counterfactual
prediction. This does not re-derive contextuality theorems; it provides a unifying structural
interpretation in which contextuality appears as a special case of reduction-induced
unfaithfulness.

\subsection{Process tensors, quantum combs, and the ``static algebra'' limitation}
\label{sec:process-tensors}

Process tensors (and closely related quantum combs) provide an operational framework for
multi-time quantum processes by representing the dependence of outcomes on inserted operations
across multiple times \cite{Pollock2018PRA,Pollock2018PRL,Chiribella2009PRA,Chiribella2008PRL}.
These objects refine $\Sigma_1$ by encoding temporal correlations, while remaining algebraically
closed: prediction proceeds by contraction rules internal to a fixed tensorial object.

The present analysis emphasizes that static algebraic closure does not, by itself, guarantee
faithfulness. If the reduction identifying a ``process object'' is degenerate with respect to
causally relevant correlations not encoded in that object, then counterfactual interventions
can distinguish histories that the $\Sigma_1$ object treats as identical. In such cases, the
process object remains a many-to-one image of $\mathcal{H}$.

Related observations appear in work showing that extension theorems for stochastic or quantum
process descriptions can fail in the presence of active interventions \cite{Milz2020Quantum}.
The $\Sigma_2$ framework isolates the minimal geometric structure required to diagnose this
failure: rather than seeking a single ``complete'' tensor, one must track continuation
structure across fiber elements or coarse-grained fiber classes.

Reference~\cite{Milz2020Quantum} shows that quantum stochastic processes need not admit
Kolmogorov extensions when interventions are present: multi-time correlations may fail to be
consistently embeddable in a single static tensor structure. This can be understood as an
instance of unfaithful reduction: a process tensor constructed from limited observational data
may fail to support counterfactual predictions under interventions not included in its
construction. The $\Sigma_2$ framework diagnoses this failure as arising from fiber degeneracy in
the map from full histories to the process tensor and characterizes the minimal fiber-aware
augmentation required to restore predictive capacity.

\subsection{Quantum causal models and process-matrix frameworks}
\label{sec:causal}

Quantum causal models and the process-matrix framework study which correlations and operational
statistics admit causal explanation, including scenarios with indefinite causal order
\cite{Oreshkov2012,CostaShrapnel2016,WoodSpekkens2015}. These approaches typically operate at the
level of correlations among nodes and constraints on admissible process objects.

The present framework addresses a prior representational question: whether the reduced
descriptions assigned to nodes support intervention semantics at all. If nodal reductions are
unfaithful with respect to the intervention class relevant to a causal analysis, then no causal
structure built atop those reduced objects can yield valid counterfactual predictions without
incorporating fiber-level admissibility structure. In this sense, $\Sigma_2$ provides a
diagnostic precondition for causal modeling in reduced operational spaces.

\subsection{Measurement and branch structure (scope-limited)}
\label{sec:measurement}

The measurement problem is often framed as a tension between unitary evolution and definite
outcomes. We make a narrower claim: many measurement puzzles can be read as instances where a
reduction $\Pi$ identifies distinct histories (e.g.\ distinct macroscopic records) under a
single reduced description, thereby destroying the continuation structure required for
counterfactual prediction about measurement outcomes.

The $\Sigma_2$ framework does not select which history in a fiber is realized. It instead
clarifies what any successful account must retain: a branch- or fiber-indexed continuation
geometry sufficient to support counterfactual intervention questions of the form
``what would occur if we measured differently?'' This establishes a structural constraint on
reduction-based resolutions without committing to a particular ontology.

\subsection{Relation to classic no-go results}
\label{sec:nogos}

Classic results such as Kochen--Specker contextuality \cite{KochenSpecker1967} and Fine's
equivalence results for Bell-type models \cite{Fine1982} constrain classes of hidden-variable
assignments compatible with quantum statistics. The present work is orthogonal in emphasis:
rather than ruling out ontological models, it constrains the \emph{form of reduced descriptive
objects} under which counterfactual prediction is even well-defined. Where those no-go results
apply, they can often be reinterpreted as special instances of unfaithful cuts relative to
particular intervention classes.